\documentclass[aspectratio=169]{beamer}
\usepackage[utf8]{inputenc}
\usepackage[T1]{fontenc}
\usepackage[slovak]{babel}
\usepackage{lmodern}
\usepackage{graphicx}
\usepackage{booktabs}

\usetheme{Madrid}
\usecolortheme{seahorse}
\setbeamertemplate{navigation symbols}{}

\title{ProMat}
\subtitle{Enterprise AI Project Management}
\author{Maturitny projekt}
\institute{SPSE Zochova 9, Bratislava}
\date{2025/2026}

\begin{document}

\begin{frame}
  \titlepage
\end{frame}

\begin{frame}{Osnova prezentacie}
  \tableofcontents
\end{frame}

\section{Motivacia a ciel}
\begin{frame}{Problem a motivacia}
  \begin{itemize}
    \item Tymy riesia ulohy, dokumenty a komunikaciu v oddelenych nastrojoch.
    \item Hladanie informacii v PDF a DOCX je pomale a neprehladne.
    \item Chyba centralny auditovatelny priestor pre rozhodovanie nad datami.
  \end{itemize}
  \vspace{0.4cm}
  \textbf{Ciel:} vytvorit lokalny system, ktory spaja projektove riadenie s AI pracou nad dokumentmi.
\end{frame}

\begin{frame}{Hlavne ciele projektu}
  \begin{itemize}
    \item Sprava projektov, uloh, clenov timu a opravneni.
    \item Nahravanie dokumentov (PDF, DOCX) a extrakcia textu.
    \item AI chat nad dokumentmi so sumarizaciou a citaciami.
    \item Dashboard pre rychly prehlad stavu projektu.
    \item Prevadzka lokalne bez odosielania dat tretim stranam.
  \end{itemize}
\end{frame}

\section{Technicke riesenie}
\begin{frame}{Technologie}
  \begin{itemize}
    \item Backend: Python, Flask, SQLAlchemy.
    \item Frontend: Jinja2, HTML/CSS, JavaScript.
    \item Databaza: SQLite.
    \item Ulozisko suborov: lokalny priecinok \texttt{uploads/}.
    \item AI vrstva: lokalny model cez Ollama.
  \end{itemize}
\end{frame}

\begin{frame}{Architektura aplikacie}
  \begin{itemize}
    \item Klient-server architektura, monoliticky Flask backend.
    \item Blueprinty rozdelene podla domen:
      \begin{itemize}
        \item autentifikacia,
        \item projekty a ulohy,
        \item dokumenty,
        \item AI chat,
        \item team a nastavenia.
      \end{itemize}
    \item Servisna vrstva pre extrakciu textu a AI logiku.
    \item Sablony a znovupouzitelne komponenty pre konzistentne UI.
  \end{itemize}
\end{frame}

\begin{frame}{Datovy model - klucove entity}
  \begin{itemize}
    \item \textbf{User}: identita, rola, pristupy.
    \item \textbf{Project}: nazov, stav, progres, vlastnik.
    \item \textbf{Task}: stav, priorita, termin, priradeny pouzivatel.
    \item \textbf{Document}: metadata suboru, extrahovany text, tagy.
    \item \textbf{ChatSession / ChatMessage}: historia AI komunikacie.
  \end{itemize}
  \vspace{0.3cm}
  Vztahy: Projekt ma ulohy a dokumenty, ulohy maju komentare, chat je naviazany na pouzivatela a volitelne na dokument.
\end{frame}

\section{Klucove funkcie}
\begin{frame}{Autentifikacia a roly}
  \begin{itemize}
    \item Registracia a prihlasenie s hashovanim hesla.
    \item Session management cez Flask-Login.
    \item Role-based pristup (admin + zakladne roly).
    \item Opravnenia riadia viditelnost projektov a manazerske akcie.
  \end{itemize}
\end{frame}

\begin{frame}{Sprava projektov a uloh}
  \begin{itemize}
    \item CRUD operacie pre projekty a ulohy.
    \item Tabulkovy aj kanban pohlad na ulohy.
    \item Zmena stavu drag-and-drop s okamzitou aktualizaciou.
    \item Filtre: projekt, stav, priorita, zodpovedna osoba.
    \item Komentare, checklisty, prilohy a aktivita pri ulohach.
  \end{itemize}
\end{frame}

\begin{frame}{Dokumenty a AI chat}
  \begin{itemize}
    \item Upload dokumentu a automaticka extrakcia textu.
    \item Verziovanie a tagovanie dokumentov.
    \item AI chat nad obsahom dokumentov.
    \item Sumarizacia a vyhladavanie pojmov.
    \item Odpovede s citaciami pre overitelnost.
  \end{itemize}
\end{frame}

\begin{frame}{Dashboard a prinosy}
  \begin{itemize}
    \item Statistiky: aktivne projekty, terminy, rozpracovane ulohy.
    \item Rychly prehlad nedavnych aktivit a dokumentov.
    \item Centralizacia dat zmensuje straty casu pri hladani.
    \item Transparentnost procesu a jednoduchsia spolupraca timu.
  \end{itemize}
\end{frame}

\section{Zaver}
\begin{frame}{Zhrnutie}
  \begin{itemize}
    \item Projekt ProMat uspesne prepaja riadenie projektov s AI pracou nad dokumentmi.
    \item Aplikacia pokryva klucove procesy: planovanie, dokumentaciu, komunikaciu a reporting.
    \item Najvacsia hodnota je v rychlom ziskavani relevantnych informacii z dokumentov.
    \item Riesene lokalne prostredie zvysuje kontrolu nad citlivymi datami.
  \end{itemize}
  \vspace{0.4cm}
  \textbf{Mozne dalsie kroky:}
  \begin{itemize}
    \item pokrocile AI metriky kvality odpovedi,
    \item notifikacie v realnom case,
    \item export reportov pre vedenie.
  \end{itemize}
\end{frame}

\begin{frame}{Dakujem za pozornost}
  Otazky?
\end{frame}

\end{document}
